% GENERATED FILE. Edit canonical lessons, then run npm run generate:books.
% canonical-source-hash: fnv1a64:56deb0f143569aa2
% canonical-lessons: HI-C117-mausi, HI-C117-bua, HI-C117-mama, HI-C117-bhabhi, HI-C117-sasural

\chapter{Aunts, Uncles, In-Laws}
\label{ch:hi-aunts-uncles}
\hlchaptermodality{117}

\begin{chapteropening}
\textbf{By the end of this chapter:} \emph{I can name aunts and uncles on both sides, a brother's wife, and the home of the in-laws.}

Complete HI-C117-sasural: the rest of the family, named by side.
\end{chapteropening}

\section[\texorpdfstring{\hi{मौसी} (mausī)}{मौसी (mausī)}]{\hi{मौसी} (mausī) — aunt, the mother's sister}

\label{lesson:HI-C117-mausi}



\begin{quote}
Before the new one: say the Hindi for the mother's mother, then the Hindi for the father's younger brother.
\end{quote}

\subsection*{You'll want to know: \hi{मौसी}}
\textbf{\hi{मौसी}} (\emph{mausī}) — \textquotedblleft{}aunt\textquotedblright{}, your mother's sister.

It is often explained as \textbf{\hi{माँ} \hi{जैसी}}, \textquotedblleft{}like a mother\textquotedblright{}: the aunt closest to a mother's place.

\begin{cousinweb}
\textbf{\hi{मौसी}} (\emph{mausī}) is a tadbhava, a word inherited from Sanskrit and reshaped by sound change. It comes from \emph{mātuḥṣvasṛ}, \textquotedblleft{}mother's sister\textquotedblright{} (CDIAL 10001). It is built from \emph{mātṛ}, \textquotedblleft{}mother\textquotedblright{}, and \emph{svasṛ}, \textquotedblleft{}sister\textquotedblright{}.
\end{cousinweb}

\begin{grammarlens}[title={What you've built}]
An aunt, on your mother's side.
\end{grammarlens}

\subsection*{Your turn}
\begin{itemize}
\raggedright
  \item \emph{Say it:} \emph{mausī}
  \item \emph{Say it:} \emph{mausī}, once more
  \item \emph{Say it:} say \emph{merī mausī}
  \item \emph{Recall:} say the Hindi for the mother's mother, then the Hindi for the father's younger brother, then say \emph{mausī} again
\end{itemize}

\begin{tcolorbox}[breakable,colback=teal!4,colframe=teal!35!black,title={Before you move on}]
What does \hi{मौसी} mean? (\textquotedblleft{}Aunt, the mother's sister\textquotedblright{}.) Say it once more.
\end{tcolorbox}
\section[\texorpdfstring{\hi{बुआ} (buā)}{बुआ (buā)}]{\hi{बुआ} (buā) — aunt, the father's sister}

\label{lesson:HI-C117-bua}



\begin{quote}
Before the new one: say the Hindi for the father's younger brother, then the Hindi for the mother's sister.
\end{quote}

\subsection*{You'll want to know: \hi{बुआ}}
\textbf{\hi{बुआ}} (\emph{buā}) — \textquotedblleft{}aunt\textquotedblright{}, your father's sister.

\textbf{\hi{मौसी}} and \textbf{\hi{बुआ}} are both \textquotedblleft{}aunt\textquotedblright{} in English; in Hindi the side of the family decides.

\begin{cousinweb}
\textbf{\hi{बुआ}} (\emph{buā}) is close to the Hindi words \emph{phūphī} and \emph{phūā}, \textquotedblleft{}father's sister\textquotedblright{}, which CDIAL lists under a reconstructed \emph{*phupphī} (CDIAL 9089). CDIAL does not list \emph{buā} itself, so its exact path is uncertain.
\end{cousinweb}

\begin{grammarlens}[title={What you've built}]
Both aunts, and the side each belongs to.
\end{grammarlens}

\subsection*{Your turn}
\begin{itemize}
\raggedright
  \item \emph{Say it:} \emph{buā}
  \item \emph{Say it:} \emph{buā}, once more
  \item \emph{Say it:} say \emph{merī mausī aur merī buā}
  \item \emph{Recall:} say the Hindi for the father's younger brother, then the Hindi for the mother's sister, then say \emph{buā} again
\end{itemize}

\begin{tcolorbox}[breakable,colback=teal!4,colframe=teal!35!black,title={Before you move on}]
What does \hi{बुआ} mean? (\textquotedblleft{}Aunt, the father's sister\textquotedblright{}.) Say it once more.
\end{tcolorbox}
\section[\texorpdfstring{\hi{मामा} (māmā)}{मामा (māmā)}]{\hi{मामा} (māmā) — uncle, the mother's brother}

\label{lesson:HI-C117-mama}



\begin{quote}
Before the new one: say the Hindi for the mother's sister, then the Hindi for the father's sister.
\end{quote}

\subsection*{You'll want to know: \hi{मामा}}
\textbf{\hi{मामा}} (\emph{māmā}) — \textquotedblleft{}uncle\textquotedblright{}, your mother's brother.

Beside \textbf{\hi{चाचा}}, the father's brother, it completes the uncles you meet most.

\begin{cousinweb}
\textbf{\hi{मामा}} (\emph{māmā}) is a tadbhava, a word inherited from Sanskrit and reshaped by sound change. It comes from \emph{māma}, \textquotedblleft{}uncle\textquotedblright{} (CDIAL 10055). CDIAL calls it a nursery word and notes similar words in Dravidian.
\end{cousinweb}

\begin{grammarlens}[title={What you've built}]
An uncle on your mother's side.
\end{grammarlens}

\subsection*{Your turn}
\begin{itemize}
\raggedright
  \item \emph{Say it:} \emph{māmā}
  \item \emph{Say it:} \emph{māmā}, once more
  \item \emph{Say it:} say \emph{mere māmā}
  \item \emph{Recall:} say the Hindi for the mother's sister, then the Hindi for the father's sister, then say \emph{māmā} again
\end{itemize}

\begin{tcolorbox}[breakable,colback=teal!4,colframe=teal!35!black,title={Before you move on}]
What does \hi{मामा} mean? (\textquotedblleft{}Uncle, the mother's brother\textquotedblright{}.) Say it once more.
\end{tcolorbox}
\section[\texorpdfstring{\hi{भाभी} (bhābhī)}{भाभी (bhābhī)}]{\hi{भाभी} (bhābhī) — sister-in-law, a brother's wife}

\label{lesson:HI-C117-bhabhi}



\begin{quote}
Before the new one: say the Hindi for the father's sister, then the Hindi for the mother's brother.
\end{quote}

\subsection*{You'll want to know: \hi{भाभी}}
\textbf{\hi{भाभी}} (\emph{bhābhī}) — \textquotedblleft{}sister-in-law\textquotedblright{}, specifically your brother's wife.

It shares its start with \textbf{\hi{भाई}}, \textquotedblleft{}brother\textquotedblright{}: she is named through him.

\begin{cousinweb}
\textbf{\hi{भाभी}} (\emph{bhābhī}) comes from a reconstructed \emph{*bhābbha-}, which CDIAL calls a nursery word; it lists Hindi \emph{bhābhī}, \textquotedblleft{}elder brother's wife\textquotedblright{} (CDIAL 9453a).
\end{cousinweb}

\begin{grammarlens}[title={What you've built}]
A relative by marriage, named through your brother.
\end{grammarlens}

\subsection*{Your turn}
\begin{itemize}
\raggedright
  \item \emph{Say it:} \emph{bhābhī}
  \item \emph{Say it:} \emph{bhābhī}, once more
  \item \emph{Say it:} say \emph{merī bhābhī}
  \item \emph{Recall:} say the Hindi for the father's sister, then the Hindi for the mother's brother, then say \emph{bhābhī} again
\end{itemize}

\begin{tcolorbox}[breakable,colback=teal!4,colframe=teal!35!black,title={Before you move on}]
What does \hi{भाभी} mean? (\textquotedblleft{}Sister-in-law, a brother's wife\textquotedblright{}.) Say it once more.
\end{tcolorbox}
\section[\texorpdfstring{\hi{ससुराल} (sasurāl)}{ससुराल (sasurāl)}]{\hi{ससुराल} (sasurāl) — the in-laws' home}

\label{lesson:HI-C117-sasural}



\begin{quote}
Before the new one: say the Hindi for the mother's brother, then the Hindi for a brother's wife.
\end{quote}

\subsection*{You'll want to know: \hi{ससुराल}}
\textbf{\hi{ससुराल}} (\emph{sasurāl}) — \textquotedblleft{}the in-laws' home\textquotedblright{}, the family a person marries into.

\textbf{\hi{वह} \hi{ससुराल} \hi{गई} \hi{है}} — \textquotedblleft{}she has gone to her in-laws' home\textquotedblright{}.

\begin{cousinweb}
\textbf{\hi{ससुराल}} (\emph{sasurāl}) is built on \emph{sasur}, \textquotedblleft{}father-in-law\textquotedblright{}, which comes from Sanskrit \emph{śvaśura} (CDIAL 12753), with the place ending \emph{-āl}.
\end{cousinweb}

\begin{grammarlens}[title={What you've built}]
Aunts, uncles, a sister-in-law, and the home of the in-laws.
\end{grammarlens}

\subsection*{Your turn}
\begin{itemize}
\raggedright
  \item \emph{Say it:} \emph{sasurāl}
  \item \emph{Say it:} \emph{sasurāl}, once more
  \item \emph{Say it:} say \emph{sasurāl}
  \item \emph{Recall:} say the Hindi for the mother's brother, then the Hindi for a brother's wife, then say \emph{sasurāl} again
\end{itemize}

\begin{tcolorbox}[breakable,colback=teal!4,colframe=teal!35!black,title={Before you move on}]
What does \hi{ससुराल} mean? (\textquotedblleft{}The in-laws' home\textquotedblright{}.) Say it once more.
\end{tcolorbox}
